
\section*{4. Phenomena specific to characteristic \(p>0\)}
\phantomsection
\label{sga3:VIIB:4}
\addcontentsline{toc}{section}{4. Phenomena specific to characteristic p > 0}

Throughout \SGARefLink{sga3:VIIB:4}{Section~4}, \(p\) denotes a prime number, \(k\) a
pseudocompact ring of characteristic \(p\), and
\[
  \pi:k\longrightarrow k,\qquad x\longmapsto x^p,
\]
the continuous Frobenius endomorphism of \(k\).

\subsection*{4.1}
\phantomsection
\label{sga3:VIIB:4.1}

Let \(X\) be a formal \(k\)-variety with affine algebra \(A\).
We write \(X^{(p/k)}\), or simply \(X^{(p)}\), for the formal
\(k\)-variety
\[
  X\widehat\otimes_{\pi}k
\]
obtained by the base change \(\pi:k\to k\)
\textup{(cf.\ \SGARefLink{sga3:VIIB:1.2.D}{1.2.D})}. Its affine algebra is the completed tensor
product
\[
  A\widehat\otimes_{\pi}k.
\]
The continuous morphism
\[
  A\widehat\otimes_{\pi}k\longrightarrow A,
  \qquad
  a\widehat\otimes_{\pi}\lambda
  \longmapsto
  a^p\lambda,
\]
therefore induces a morphism of formal \(k\)-varieties
\[
  \operatorname{Fr}(X/k):
  X\longrightarrow X^{(p/k)}.
\]
We call \(\operatorname{Fr}(X/k)\) the \emph{Frobenius morphism of
\(X\) relative to \(k\)}, and below we shall often write
\(\operatorname{Fr}\) in place of \(\operatorname{Fr}(X/k)\).

